% unidades/20-clausulas-relativas.tex
\chapter{🔗 \ruby{関係節}{かんけいせつ} — Cláusulas Relativas Avanzadas}
\unidadheader{20}{関係節}{複雑な関係節}{Modificar sustantivos con cláusulas verbales complejas}

\begin{tcolorbox}[vocabbox, title={\emoji{pushpin} Vocabulario}]
\begin{tabularx}{\linewidth}{llX}
\toprule
\textbf{Japonés} & \textbf{Español} & \textbf{Tipo} \\ \midrule
\vocabitem{事件}{じけん}{incidente / caso}{sustantivo}
\vocabitem{原因}{げんいん}{causa / razón}{sustantivo}
\vocabitem{場所}{ばしょ}{lugar}{sustantivo}
\vocabitem{理由}{りゆう}{razón / motivo}{sustantivo}
\vocabitem{方法}{ほうほう}{método / manera}{sustantivo}
\vocabitem{人物}{じんぶつ}{persona / personaje}{sustantivo}
\vocabitem{記念日}{きねんび}{aniversario / día especial}{sustantivo}
\bottomrule
\end{tabularx}
\end{tcolorbox}

\subsection*{\emoji{speech-balloon} Frases Clave}
\frase{\ruby{私}{わたし}が\ruby{読}{よ}んだ\ruby{本}{ほん}はとても\ruby{面白}{おもしろ}かった。}{El libro que leí fue muy interesante.}
\frase{\ruby{彼女}{かのじょ}が\ruby{作}{つく}ったケーキを\ruby{食}{た}べた。}{Comí el pastel que ella hizo.}
\frase{\ruby{田中}{たなか}さんに\ruby{会}{あ}った\ruby{場所}{ばしょ}に\ruby{行}{い}きたい。}{Quiero ir al lugar donde conocí a Tanaka.}
\frase{\ruby{子供}{こども}が\ruby{生}{う}まれた\ruby{日}{ひ}を\ruby{忘}{わす}れられない。}{No puedo olvidar el día que nació mi hijo.}

\subsection*{\emoji{gear} Gramática}
\patron{Cláusula relativa simple — modificar sustantivos}
  {[cláusula verbal en plain form] + sustantivo}
  {\ej{\ruby{昨日}{きのう}\ruby{読}{よ}んだ\ruby{本}{ほん}。}{El libro que leí ayer.}
   \ej{\ruby{私}{わたし}が\ruby{住}{す}んでいる\ruby{町}{まち}。}{El pueblo donde vivo.}
   \ej{\ruby{田中}{たなか}さんが\ruby{作}{つく}った\ruby{料理}{りょうり}。}{La comida que hizo Tanaka.}}

\patron{Cláusulas relativas con sustantivos de lugar / tiempo / razón}
  {lugar: [cláusula] + \ruby{場所}{ばしょ}/\ruby{ところ} \quad tiempo: [cláusula] + \ruby{日}{ひ}/\ruby{時}{とき}}
  {\ej{\ruby{彼}{かれ}と\ruby{初}{はじ}めて\ruby{会}{あ}った\ruby{場所}{ばしょ}に\ruby{行}{い}った。}{Fui al lugar donde lo conocí por primera vez.}
   \ej{\ruby{子供}{こども}が\ruby{生}{う}まれた\ruby{日}{ひ}は\ruby{忘}{わす}れない。}{No olvidaré el día en que nació mi hijo.}
   \ej{\ruby{卒業}{そつぎょう}した\ruby{年}{とし}に\ruby{旅行}{りょこう}した。}{Viajé el año que me gradué.}}

\patron{Nominalización de cláusula con の/こと como sujeto}
  {[cláusula plain] + の/こと + が/は/を}
  {\ej{\ruby{彼}{かれ}が\ruby{来}{こ}なかったのは\ruby{残念}{ざんねん}だ。}{Es una lástima que él no haya venido.}
   \ej{\ruby{彼女}{かのじょ}が\ruby{合格}{ごうかく}したことを\ruby{知}{し}っていた。}{Sabía que ella había aprobado.}
   \ej{\ruby{雨}{あめ}が\ruby{降}{ふ}っているのが\ruby{見}{み}える。}{Se puede ver que está lloviendo.}}

\comparacion{Cláusula relativa vs cláusula nominalizada}
  {Cláusula relativa\\\small modifica sustantivo explícito\\\small \ruby{読}{よ}んだ\ruby{本}{ほん}\\\textit{el libro que leí}}
  {Cláusula nominalizada\\\small el verbo ES el sustantivo\\\small \ruby{読}{よ}んだのが\ruby{好}{す}き\\\textit{lo que leí me gusta}}
  {En la cláusula relativa, hay un sustantivo cabeza (本). En la nominalizada, の o こと reemplaza ese sustantivo y la cláusula completa funciona como sujeto/objeto.}

\coloquial{SOV y cláusulas relativas}{
  En casual, las cláusulas relativas son idénticas. El reto es la longitud: en japonés, todos los modificadores van ANTES del sustantivo. «El libro que encontré ayer en la librería cerca de la estación» → \ruby{駅}{えき}\ruby{近}{ちか}くの\ruby{本屋}{ほんや}で\ruby{昨日}{きのう}\ruby{見}{み}つけた\ruby{本}{ほん}. ¡El modificador puede ser muy largo!}

\culturabox{SOV\ruby{構造}{こうぞう} — El orden de palabras japonés}{
  El japonés es SOV (Sujeto-Objeto-Verbo) y todas las cláusulas modificadoras van ANTES del sustantivo que modifican. Esto es lo opuesto al español, donde las cláusulas relativas van después (el libro QUE leí). Para los hispanohablantes, leer japonés implica aprender a «leer al revés»: primero el modificador, luego el sustantivo.
}

\lectura{\ruby{思}{おも}い\ruby{出}{で}の\ruby{場所}{ばしょ}\\[0.4em]
  \ruby{私}{わたし}が\ruby{生}{う}まれた\ruby{町}{まち}はとても\ruby{小}{ちい}さい。\ruby{子供}{こども}のとき\ruby{遊}{あそ}んだ\ruby{公園}{こうえん}はまだある。\ruby{初}{はじ}めて\ruby{彼女}{かのじょ}に\ruby{会}{あ}った\ruby{場所}{ばしょ}は\ruby{閉}{し}まってしまった。\ruby{昔}{むかし}\ruby{住}{す}んでいた\ruby{家}{いえ}は\ruby{新}{あたら}しい\ruby{建物}{たてもの}になっていた。\ruby{変}{か}わらないのは\ruby{山}{やま}の\ruby{景色}{けしき}だけだ。}
{El lugar de los recuerdos\\[0.4em]
  El pueblo donde nací es muy pequeño. El parque donde jugé de niño todavía existe. El lugar donde conocí a mi novia por primera vez ha cerrado. La casa donde vivía antes se ha convertido en un edificio nuevo. Lo único que no ha cambiado es el paisaje de la montaña.}

\begin{tcolorbox}[ejerciciobox, title={\emoji{pencil} Ejercicios}]
\begin{enumerate}
  \item Combina: \ruby{本}{ほん} / \ruby{昨日}{きのう}\ruby{買}{か}った → \_\_\_
  \item Construye con の: \ruby{田中}{たなか}さんが\ruby{来}{こ}なかった + \ruby{残念}{ざんねん}だ
  \item Describe tu ciudad usando cláusulas relativas: \ruby{生}{う}まれた\ruby{場所}{ばしょ}は...
\end{enumerate}
\textbf{Respuestas:}
1) \ruby{昨日}{きのう}\ruby{買}{か}った\ruby{本}{ほん} \quad
2) \ruby{田中}{たなか}さんが\ruby{来}{こ}なかったのは\ruby{残念}{ざんねん}だ \quad
3) (libre)
\end{tcolorbox}

\begin{tcolorbox}[kanjibox, title={\emoji{mahjong} Kanjis}]
\begin{tabularx}{\linewidth}{cllXX}
\toprule
\textbf{Kanji} & \textbf{On} & \textbf{Kun} & \textbf{Significado} & \textbf{Vocab} \\ \midrule
\kanjirow{原}{げん}{\ruby{もと}{}}{origen/causa}{\ruby{原因}{げんいん} / \ruby{原}{もと}}
\kanjirow{因}{いん}{\ruby{よ}{る}}{causa/basarse}{\ruby{原因}{げんいん} / \ruby{因果}{いんが}}
\kanjirow{理}{り}{—}{razón/lógica}{\ruby{理由}{りゆう} / \ruby{道理}{どうり}}
\kanjirow{由}{ゆう}{\ruby{よし}{}}{razón/fuente}{\ruby{理由}{りゆう} / \ruby{由来}{ゆらい}}
\bottomrule
\end{tabularx}
\end{tcolorbox}
